\section{The coupled adjoint equations}\label{sec:deformation}
Choose one reference Kummer representation over $\Oloc/49$,
\begin{equation}\label{eq:rho0}
\rho_0(g)=\begin{pmatrix}\chi(g)&b_0(g)\\0&1\end{pmatrix},
\end{equation}
whose reduction has the same residual $\alpha$ as the actual natural representation. The reference is built by lifting the coefficients of the same global classes $[\phi],[3],[\sqrt5]$. It has no seven valuation coordinate at the fixed place. It is an auxiliary representation, not the actual one.

Every determinant $\chi$ framed lift of this residual representation has a unique expression
\begin{equation}\label{eq:adjoint}
\rho(g)=(I+7X(g))\rho_0(g),\qquad
X(g)=\begin{pmatrix}p(g)&q(g)\\\beta(g)&-p(g)\end{pmatrix}\in\mathfrak{sl}_2(k).
\end{equation}
The determinant condition makes $X$ traceless. Multiplication shows that $\rho$ is a representation if and only if
\[
X(gh)=X(g)+\Ad(\overline\rho_0(g))X(h).
\]
For $\overline\rho_0(g)=\left(\begin{smallmatrix}x&a\\0&1\end{smallmatrix}\right)$, direct conjugation sends $(p,q,b)$ to
\[
(p+ax^{-1}b,\;xq-2ap-a^2x^{-1}b,\;x^{-1}b).
\]
Thus the three coupled equations are
\begin{align}
\beta(gh)&=\beta(g)+\chi(g)^{-1}\beta(h),\label{eq:beta-cocycle}\\
p(gh)&=p(g)+p(h)+\alpha(g)\chi(g)^{-1}\beta(h),\label{eq:p-cocycle}\\
q(gh)&=q(g)+\chi(g)q(h)-2\alpha(g)p(h)
-\alpha(g)^2\chi(g)^{-1}\beta(h).\label{eq:q-cocycle}
\end{align}
All coefficients in these equations lie in $k$. The reference itself is a representation over $\Oloc/49$, so no characteristic seven carry term has been silently discarded in deriving them.

Throughout, $\alpha$ and beta are normalized continuous cocycles with their displayed coefficient actions. Their cup has trivial target action and is explicitly $(\alpha\cup\beta)(g,h)=\alpha(g)\chi(g)^{-1}\beta(h)$. Define the ordinary differential by
\[
dp(g,h)=p(g)+p(h)-p(gh)
\]
and the twisted differential by $d_\chi q(g,h)=q(g)+\chi(g)q(h)-q(gh)$. Then
\begin{align}
dp&=-\alpha\cup\beta,\label{eq:p-equation}\\
d_\chi q(g,h)&=2\alpha(g)p(h)+\alpha(g)^2\chi(g)^{-1}\beta(h).\label{eq:q-equation}
\end{align}
The right side of \eqref{eq:q-equation} is a two cocycle once the preceding equations hold. For a general fixed $p$, its vanishing in cohomology is the next bare existence problem. Solving the $p$ equation alone does not formally solve the $q$ equation. The project does not establish that this bare upper class is nonzero in the actual branches. Prescribed upper geometry and integral lifting retain additional conditions.

The actual trace is
\begin{equation}\label{eq:adjointtrace}
\tr\rho(g)=\chi(g)+1
+7\bigl((\chi(g)-1)p(g)+\alpha(g)\beta(g)\bigr)\pmod{49}.
\end{equation}
The upper cochain $q$ is invisible to trace at this precision. On $\chi(g)=1$, the diagonal cochain $p$ is also invisible and only the product $\alpha\beta$ remains. This explains why a homogeneous Frobenius sieve can leave essential deformation coordinates undetermined.

Strict conjugation by $I+7Y$ changes $X$ by $Y-\Ad(\overline\rho_0)Y$. Independent local changes can simplify representatives but do not create a common global integral frame. The affine target in the last part of this reader is formulated after one frame and one reference have been fixed.

\subsection{A delayed tame obstruction}
Suppose at a beta support prime the actual inertia and Frobenius have
\[
I+7M\pmod{49},\qquad I+7A\pmod{49},\qquad Q=1+7h\pmod{49}.
\]
The tame relation is automatically insensitive to some first digit changes modulo forty nine. One precision later it forces
\begin{equation}\label{eq:tame-next}
[A,M]=hM\quad\text{over }k.
\end{equation}
A trace invisible correction modulo forty nine need not pass \eqref{eq:tame-next}. The archive retains a seven valuation correction that fails this next relation at $883$ when complete inertia is held fixed. This excludes that strengthening inside the stipulated deformation fiber. It does not prove that the actual valuation coordinate is zero in every candidate, because equality of a finite trace list is weaker than equality of the entire framed inertia and trace data used by the fiber theorem.
